\section{Discussão e ameaças à validade}

\textbf{Síntese das hipóteses.} H1 é sustentada; H2, H4, H6 e H7, parcialmente;
H3 e H5 não são sustentadas. A frequência típica é baixa, enquanto recuperação de
CI é rápida. Comunidades maiores associam-se a lead times maiores, contrariando
uma interpretação simplista de que mais contribuidores garantem velocidade.
A divergência mais importante é de construto: trocar falhas de CI por releases
corretivas altera tanto a relação com frequência quanto a classificação geral.
É plausível que publicar frequentemente aumente a oportunidade de observar outra
release em sete dias, mas o estudo não identifica esse mecanismo causalmente.

\textbf{Construto.} Releases e tags não equivalem a deploys; CI não é produção.
Tags usam autoria do commit, suscetível a rebase/squash, e podem não representar
a criação da versão. O CFR~(b) depende de heurística imperfeita, de versões
paralelas e da janela de sete dias. Rework também não demonstra retrabalho causado
por incidente. Faixas fixas da disciplina não são benchmarks universais atuais.
O arredondamento da mediana afeta \labACOneFracionarias\ classificações de C1:
notas centrais diferentes deixam projetos próximos a um corte.

\textbf{Interna e de conclusão.} Idade, domínio, popularidade, governança e tamanho
dos projetos não foram controlados causalmente. Compartilhar ecossistema pode
introduzir dependência entre repositórios. Subgrupos desiguais (tipo outro tem só
seis recuperações disponíveis), ausências informativas e muitos testes limitam
inferências. Reportamos n, efeitos, intervalos e Holm, preservando testes não
significativos e resultados não estimáveis. As famílias exploratórias maiores
alteram a significância de tipo--recuperação, embora os dados não tenham mudado.

\textbf{Validação e validade externa.} Três avaliações humanas foram declaradas
pelos integrantes; hashes verificam integridade, não independência cognitiva.
Os 240 tipos complementares tiveram uma avaliação por projeto após sugestões
automáticas, com possível ancoragem. Refinar e avaliar a v3 na mesma amostra pode
superestimar F1; não houve teste separado de generalização. Seleção por popularidade,
uso de Actions e mínimo de releases/runs exclui projetos menores, inativos,
corporativos e outros sistemas de CI. A janela de um ano não permite extrapolar
para toda a história dos projetos. Reprodutibilidade offline não elimina essas
limitações e não equivale à replicação independente por outro grupo.

\section{Conclusão e material de reprodução}

Os proxies permitem estudar a cadência e a recuperação de CI, mas conclusões sobre
estabilidade dependem da definição. As categorias mudam em cerca de 40\% dos pares
C1--C2/C3, apesar de os dados serem os mesmos. Recomendamos reportar variantes,
ausências e denominadores, evitando rótulos de desempenho sem explicitar os proxies.
O bônus complementa CFR~(b) ao contar corretivas, não as versões que as precedem.

Código, dataset congelado, fontes, testes, notebooks, dicionários e comandos de
reprodução estão no repositório do grupo:
\begin{flushleft}\footnotesize
\url{https://github.com/PedroMaiaAlves/lab-experimentacao-software}
\end{flushleft}
O índice de projetos associado ao repositório está em
\url{https://github.com/PedroMaiaAlves/lab-experimentacao-software/projects}.
A execução local não comprova publicação na main, atualização do quadro ou
replicação cruzada. Sem repositório sorteado, dataset comparável e registros de
execução externa, não declaramos essa etapa concluída nem inventamos Issues ou
respostas de outros grupos.
